\documentclass{article}
\usepackage{caption}
\usepackage{booktabs}
\usepackage{amssymb}
\usepackage{graphicx}
\usepackage{pdflscape}
\usepackage{float}
\usepackage[margin=1in]{geometry}

% For esttab significance stars like \sym{***}
\newcommand{\sym}[1]{\ifmmode^{#1}\else\(^{#1}\)\fi}

\begin{document}


\begin{table}[!htbp]
\centering
\caption*{Table 1. Summary Statistics for Plot-Level Analysis}
\begin{tabular}{llllll}
    \midrule
                    &         Obs&        Mean&   Std. Dev.&         Min&         Max\\
\midrule
Wildfire occurrence &      761250&        .055&        .227&           0&           1\\
Harvest             &      761250&        .011&        .103&           0&           1\\
Thinning            &      761250&        .006&         .08&           0&           1\\
Harvest shift-share &      761250&        .013&        .009&           0&        .172\\
Thinning shift-share&      761250&        .007&        .011&           0&        .076\\
Total forest area (thousand acre)&      761250&        5.83&       3.213&        .001&      44.892\\
Maximum VPD (hPa)   &      761250&      19.079&       4.071&        9.95&        44.3\\
Trees per acre      &      761250&     328.469&     259.092&           0&     999.662\\
Sawtimber volume per acre (ft$^3$/acre)&      761250&      699.04&    1031.797&           0&    17464.57\\
Stand age           &      761250&      40.991&      26.112&           0&         200\\
Percentage artifical forest&      761250&        .146&         .33&           0&           1\\
Slope (\%)          &      761250&       8.832&      14.438&           0&         155\\
Elevation (m)       &      761250&     752.652&     800.871&           0&        7170\\
Bio growth rate     &      761250&        .967&       1.648&    -146.259&      55.621\\
Forestry and logging establishments&      761250&       5.233&       7.459&           0&          55\\
Distance to nearest road (mile)&      761250&        .394&        .553&           0&           6\\
Siteclass (Best:1 to Worst:7)&      761250&       5.103&       1.294&           1&           7\\
\midrule

\end{tabular}
\begin{minipage}{\textwidth}
\footnotesize
\textit{Notes:} Summary statistics are reported for a plot-year panel covering  37071 plots observed from 2000 to 2024, summing up to 761,250 observations. 
Stand age is capped at 200 years in the FIA data; 5\% of plots have an average stand age greater than 82 years, and 1\% have an average stand age greater than 101 years.
Bio growth rate is the annual aboveground biomass growth rate of a stem (short tons/ year). 
Site productivity captures the inherent ability of forest land to support timber growth, defined by potential 
wood production (ft$^3$/acre/year) under fully stocked conditions.
\end{minipage}
\end{table}


\begin{table}[!htbp]
\centering
\caption*{Table 2. Summary Statistics for County-Level Analysis}
\resizebox{\textwidth}{!}{\begin{tabular}{lllll}
    \midrule
                    &        Mean&   Std. Dev.&         Min&         Max\\
\midrule
Wildfire area in MTBS (thousand acre)&        .753&       6.972&           0&     348.916\\
Wildfire occurrence in MTBS&        .092&        .289&           0&           1\\
Harvest area (thousand acre)&       2.449&       5.179&           0&      70.224\\
Thinning area (thousand acre)&       1.472&       4.024&           0&      47.844\\
Harvest shift-share (thousand acre)&       2.564&       2.759&           0&      42.985\\
Thinning shift-share (thousand acre)&       1.524&        2.36&           0&        23.3\\
Maximum VPD (hPa)   &      31.774&       9.246&      14.311&      78.167\\
Trees per acre      &     510.997&     225.387&           0&     998.441\\
Total forest area (thousand acre)&      209.12&     148.844&        .549&    1637.981\\
Sawtimber volume per acre (ft$^3$/acre)&     870.075&     682.519&           0&    11163.57\\
Bio growth rate     &        1.31&        .951&       -5.66&       7.321\\
Slope (\%)          &      10.143&      11.665&           0&          65\\
Elevation (m)       &     798.063&     807.996&           0&    5487.093\\
Stand age           &      42.738&      13.793&           0&         135\\
Site productivity   &      66.975&      31.648&         9.5&         225\\
Forestry and logging establishments&       3.783&       6.055&           0&          55\\
\midrule

\end{tabular}}
\begin{minipage}{\textwidth}
\footnotesize
\textit{Notes:} Summary statistics are reported for a balanced county-year panel covering 1,263 counties observed annually from 2000 to 2023.
 Bio growth rate is the annual aboveground biomass growth rate of a stem (short tons/ year). 
Site productivity captures the inherent ability of forest land to support timber growth, defined by potential 
wood production (ft$^3$/acre/year) under fully stocked conditions. 73 counties excluded because of unrealistic trees per acre information.
\end{minipage}
\end{table}





\begin{table}[!htbp]
\centering
\caption*{Table 3.  Direct mitigation benefits from fuel reduction treatment on wildfire risk}
\resizebox{\textwidth}{!}{\begin{tabular}{llllll}
\begin{tabular}{lccccc}
\toprule
& \multicolumn{2}{c}{Harvest} & \multicolumn{2}{c}{Thinning} & \multicolumn{1}{c}{Joint IV} \\
\cmidrule(lr){2-3}\cmidrule(lr){4-5}\cmidrule(lr){6-6}
& (1) & (2) & (3) & (4) & (5) \\
& \multicolumn{1}{c}{Wildfire event} & \multicolumn{1}{c}{Wildfire event} & \multicolumn{1}{c}{Wildfire event} & \multicolumn{1}{c}{Wildfire event} & \multicolumn{1}{c}{Wildfire event} \\
& FE & IV & FE & IV & IV \\
\midrule
Harvest             &     -0.0000         &     -0.4453\sym{***}&                     &                     &     -0.3883\sym{***}\\
                    &    (0.0014)         &    (0.0830)         &                     &                     &    (0.0824)         \\
\addlinespace
Thinning            &                     &                     &     -0.0014         &     -0.1872\sym{***}&     -0.1663\sym{***}\\
                    &                     &                     &    (0.0018)         &    (0.0435)         &    (0.0471)         \\
\midrule
Kleibergen-Paap Wald F statistic&                     &      214.64         &                     &      488.11         &      107.14         \\
Observations        &     761,236         &     761,236         &     761,236         &     761,236         &     761,236         \\
Plot FE             &  \checkmark         &  \checkmark         &  \checkmark         &  \checkmark         &  \checkmark         \\
Year FE             &  \checkmark         &  \checkmark         &  \checkmark         &  \checkmark         &  \checkmark         \\
Species group FE    &  \checkmark         &  \checkmark         &  \checkmark         &  \checkmark         &  \checkmark         \\
Ecoregion FE        &  \checkmark         &  \checkmark         &  \checkmark         &  \checkmark         &  \checkmark         \\
\bottomrule
\end{tabular}

\end{tabular}}
\begin{minipage}{\textwidth}
\footnotesize
*** $p<0.01$, ** $p<0.05$, * $p<0.1$.\\
\textit{Notes:} The differentce between harvest and thinning is statistically significant with p-value=0.034.
14 singleton observations are dropped in the IV models because they do not provide within-group variation after absorbing the fixed effects.
\end{minipage}
\end{table}





\begin{table}[!htbp]
\centering
\caption*{Table 4. Combined mitigation benefits from fuel reduction treatment on wildfire risk at county level}

\resizebox{\textwidth}{!}{%
\resizebox{\textwidth}{!}{\begin{tabular}{l*{5}{c}}\midrule
                    &         Obs&        Mean&   Std. Dev.&         Min&         Max\\
\midrule
Wildfire occurrence &      761250&        .055&        .227&           0&           1\\
Harvest             &      761250&        .011&        .103&           0&           1\\
Thinning            &      761250&        .006&         .08&           0&           1\\
Harvest shift-share &      761250&        .013&        .009&           0&        .172\\
Thinning shift-share&      761250&        .007&        .011&           0&        .076\\
Total forest area (thousand acre)&      761250&        5.83&       3.213&        .001&      44.892\\
Maximum VPD (hPa)   &      761250&      19.079&       4.071&        9.95&        44.3\\
Trees per acre      &      761250&     328.469&     259.092&           0&     999.662\\
Sawtimber volume per acre (ft$^3$/acre)&      761250&      699.04&    1031.797&           0&    17464.57\\
Stand age           &      761250&      40.991&      26.112&           0&         200\\
Percentage artifical forest&      761250&        .146&         .33&           0&           1\\
Slope (\%)          &      761250&       8.832&      14.438&           0&         155\\
Elevation (m)       &      761250&     752.652&     800.871&           0&        7170\\
Bio growth rate     &      761250&        .967&       1.648&    -146.259&      55.621\\
Forestry and logging establishments&      761250&       5.233&       7.459&           0&          55\\
Distance to nearest road (mile)&      761250&        .394&        .553&           0&           6\\
Siteclass (Best:1 to Worst:7)&      761250&       5.103&       1.294&           1&           7\\
\midrule

\end{tabular}%
}}

\vspace{0.5em}

\begin{minipage}{\textwidth}
\footnotesize
*** $p<0.01$, ** $p<0.05$, * $p<0.1$.\\
\textit{Notes:} The differentce between harvest and thinning is not statistically significant with p-value= 0.600. 

%  ( 1)  logavg_harvest - logavg_thin = 0



% ------------------------------------------------------------------------------
% logmtbs_acre |      Coef.   Std. Err.      t    P>|t|     [95% Conf. Interval]
% -------------+----------------------------------------------------------------
%          (1) |  -.0853609   .1617271    -0.53   0.598    -.4026446    .2319228
% ------------------------------------------------------------------------------





\end{minipage}
\end{table}



\begin{table}[!htbp]
\centering
\caption*{Table A1. Direct effects of fuel reduction treatment on wildfire risk}
\resizebox{\textwidth}{!}{\begin{tabular}{lllll}
\begin{tabular}{lcccccccc}
\toprule
& \multicolumn{3}{c}{Harvest} & \multicolumn{3}{c}{Thinning} & \multicolumn{1}{c}{Joint IV} & \multicolumn{1}{c}{Thinning vs. Wait} \\
\cmidrule(lr){2-4}\cmidrule(lr){5-7}\cmidrule(lr){8-8}\cmidrule(lr){9-9}
& (1) & (2) & (3) & (4) & (5) & (6) & (7) & (8) \\
& \multicolumn{1}{c}{Wildfire event} & \multicolumn{1}{c}{Harvest event} & \multicolumn{1}{c}{Wildfire event} & \multicolumn{1}{c}{Wildfire event} & \multicolumn{1}{c}{Thinning event} & \multicolumn{1}{c}{Wildfire event} & \multicolumn{1}{c}{Wildfire event} & \multicolumn{1}{c}{Wildfire event} \\
& FE & First stage & IV & FE & First stage & IV & IV & IV \\
\midrule
Harvest             &     -0.0000         &                     &     -0.4453\sym{***}&                     &                     &                     &     -0.3883\sym{***}&                     \\
                    &    (0.0014)         &                     &    (0.0830)         &                     &                     &                     &    (0.0824)         &                     \\
\addlinespace
Thinning            &                     &                     &                     &     -0.0014         &                     &     -0.1872\sym{***}&     -0.1663\sym{***}&      0.2166\sym{**} \\
                    &                     &                     &                     &    (0.0018)         &                     &    (0.0435)         &    (0.0471)         &    (0.0978)         \\
\addlinespace
No treatment / wait &                     &                     &                     &                     &                     &                     &                     &      0.3871\sym{***}\\
                    &                     &                     &                     &                     &                     &                     &                     &    (0.0822)         \\
\addlinespace
Harvest shift-share instrument&                     &      0.6116\sym{***}&                     &                     &                     &                     &                     &                     \\
                    &                     &    (0.0417)         &                     &                     &                     &                     &                     &                     \\
\addlinespace
Thinning shift-share instrument&                     &                     &                     &                     &      0.7804\sym{***}&                     &                     &                     \\
                    &                     &                     &                     &                     &    (0.0353)         &                     &                     &                     \\
\addlinespace
Maximum VPD         &      0.0118\sym{***}&      0.0001         &      0.0118\sym{***}&      0.0118\sym{***}&      0.0001         &      0.0117\sym{***}&      0.0118\sym{***}&      0.0118\sym{***}\\
                    &    (0.0007)         &    (0.0001)         &    (0.0007)         &    (0.0007)         &    (0.0001)         &    (0.0007)         &    (0.0007)         &    (0.0007)         \\
\addlinespace
Trees per acre      &      0.0005\sym{**} &     -0.0210\sym{***}&     -0.0089\sym{***}&      0.0005\sym{**} &     -0.0037\sym{***}&     -0.0002         &     -0.0083\sym{***}&     -0.0083\sym{***}\\
                    &    (0.0002)         &    (0.0006)         &    (0.0018)         &    (0.0002)         &    (0.0004)         &    (0.0003)         &    (0.0017)         &    (0.0017)         \\
\addlinespace
Sawtimber volume per acre&     -0.0006         &     -0.0092\sym{***}&     -0.0044\sym{***}&     -0.0006         &     -0.0041\sym{***}&     -0.0013         &     -0.0046\sym{***}&     -0.0047\sym{***}\\
                    &    (0.0009)         &    (0.0016)         &    (0.0014)         &    (0.0009)         &    (0.0014)         &    (0.0010)         &    (0.0013)         &    (0.0013)         \\
\addlinespace
Stand age           &     -0.0006         &     -0.0099\sym{***}&     -0.0051\sym{***}&     -0.0006         &      0.0087\sym{***}&      0.0010         &     -0.0031\sym{***}&     -0.0031\sym{***}\\
                    &    (0.0005)         &    (0.0008)         &    (0.0010)         &    (0.0005)         &    (0.0005)         &    (0.0006)         &    (0.0011)         &    (0.0011)         \\
\addlinespace
Slope               &     -0.0001         &      0.0005\sym{***}&      0.0001         &     -0.0001         &     -0.0001         &     -0.0001         &      0.0001         &      0.0001         \\
                    &    (0.0001)         &    (0.0001)         &    (0.0001)         &    (0.0001)         &    (0.0001)         &    (0.0001)         &    (0.0001)         &    (0.0001)         \\
\addlinespace
Elevation           &      0.0044         &     -0.0148\sym{**} &     -0.0021         &      0.0044         &     -0.0007         &      0.0043         &     -0.0014         &     -0.0014         \\
                    &    (0.0059)         &    (0.0060)         &    (0.0065)         &    (0.0059)         &    (0.0006)         &    (0.0060)         &    (0.0064)         &    (0.0064)         \\
\addlinespace
Bio growth rate     &     -0.0010\sym{***}&     -0.0012\sym{***}&     -0.0017\sym{***}&     -0.0010\sym{***}&      0.0057\sym{***}&      0.0002         &     -0.0005         &     -0.0005         \\
                    &    (0.0004)         &    (0.0004)         &    (0.0004)         &    (0.0004)         &    (0.0009)         &    (0.0005)         &    (0.0005)         &    (0.0005)         \\
\addlinespace
Forestry establishments&      0.0008\sym{**} &     -0.0001         &      0.0007\sym{*}  &      0.0008\sym{**} &     -0.0002         &      0.0006\sym{*}  &      0.0006         &      0.0006         \\
                    &    (0.0004)         &    (0.0002)         &    (0.0004)         &    (0.0004)         &    (0.0002)         &    (0.0004)         &    (0.0004)         &    (0.0004)         \\
\midrule
Kleibergen-Paap Wald F statistic&                     &                     &      214.64         &                     &                     &      488.11         &      107.14         &      108.03         \\
Observations        &     761,236         &     761,236         &     761,236         &     761,236         &     761,236         &     761,236         &     761,236         &     761,236         \\
Plot FE             &  \checkmark         &  \checkmark         &  \checkmark         &  \checkmark         &  \checkmark         &  \checkmark         &  \checkmark         &  \checkmark         \\
Year FE             &  \checkmark         &  \checkmark         &  \checkmark         &  \checkmark         &  \checkmark         &  \checkmark         &  \checkmark         &  \checkmark         \\
Species group FE    &  \checkmark         &  \checkmark         &  \checkmark         &  \checkmark         &  \checkmark         &  \checkmark         &  \checkmark         &  \checkmark         \\
Ecoregion FE        &  \checkmark         &  \checkmark         &  \checkmark         &  \checkmark         &  \checkmark         &  \checkmark         &  \checkmark         &  \checkmark         \\
\bottomrule
\end{tabular}

\end{tabular}}
\begin{minipage}{\textwidth}
\footnotesize
*** $p<0.01$, ** $p<0.05$, * $p<0.1$.\\
\textit{Notes:} 14 singleton observations are dropped in all models because they do not provide within-group variation after absorbing the fixed effects.
\end{minipage}
\end{table}

\begin{table}[!htbp]
\centering
\caption*{Table A2. Robustness tests with alternative treatment windows}
\resizebox{\textwidth}{!}{\begin{tabular}{lllll}
\begin{tabular}{lccccccccc}
\toprule
& \multicolumn{3}{c}{Harvest} & \multicolumn{3}{c}{Thinning} & \multicolumn{3}{c}{Joint IV} \\
\cmidrule(lr){2-4}\cmidrule(lr){5-7}\cmidrule(lr){8-10}
& (1) & (2) & (3) & (4) & (5) & (6) & (7) & (8) & (9) \\
& Wildfire event & Wildfire event & Wildfire event & Wildfire event & Wildfire event & Wildfire event & Wildfire event & Wildfire event & Wildfire event \\
\midrule
Harvest (3-year)&  -0.4859\sym{***}&                  &                  &                  &                  &                  &  -0.4018\sym{***}&                  &                  \\
                & (0.1022)         &                  &                  &                  &                  &                  & (0.1072)         &                  &                  \\
\addlinespace
Harvest (4-year)&                  &  -0.4453\sym{***}&                  &                  &                  &                  &                  &  -0.3883\sym{***}&                  \\
                &                  & (0.0830)         &                  &                  &                  &                  &                  & (0.0824)         &                  \\
\addlinespace
Harvest (5-year)&                  &                  &  -0.4327\sym{***}&                  &                  &                  &                  &                  &  -0.3980\sym{***}\\
                &                  &                  & (0.0751)         &                  &                  &                  &                  &                  & (0.0713)         \\
\addlinespace
Thinning (3-year)&                  &                  &                  &  -0.3295\sym{***}&                  &                  &  -0.2610\sym{***}&                  &                  \\
                &                  &                  &                  & (0.0603)         &                  &                  & (0.0667)         &                  &                  \\
\addlinespace
Thinning (4-year)&                  &                  &                  &                  &  -0.1872\sym{***}&                  &                  &  -0.1663\sym{***}&                  \\
                &                  &                  &                  &                  & (0.0435)         &                  &                  & (0.0471)         &                  \\
\addlinespace
Thinning (5-year)&                  &                  &                  &                  &                  &  -0.0793\sym{**} &                  &                  &  -0.0891\sym{**} \\
                &                  &                  &                  &                  &                  & (0.0338)         &                  &                  & (0.0377)         \\
\midrule
Kleibergen-Paap Wald F statistic&   217.04         &   214.64         &   191.47         &   418.96         &   488.11         &   557.45         &    93.84         &   107.14         &   105.77         \\
Observations    &  774,295         &  761,236         &  746,245         &  774,295         &  761,236         &  746,245         &  774,295         &  761,236         &  746,245         \\
Plot FE         &\checkmark         &\checkmark         &\checkmark         &\checkmark         &\checkmark         &\checkmark         &\checkmark         &\checkmark         &\checkmark         \\
Year FE         &\checkmark         &\checkmark         &\checkmark         &\checkmark         &\checkmark         &\checkmark         &\checkmark         &\checkmark         &\checkmark         \\
Species group FE&\checkmark         &\checkmark         &\checkmark         &\checkmark         &\checkmark         &\checkmark         &\checkmark         &\checkmark         &\checkmark         \\
Ecoregion FE    &\checkmark         &\checkmark         &\checkmark         &\checkmark         &\checkmark         &\checkmark         &\checkmark         &\checkmark         &\checkmark         \\
\bottomrule
\end{tabular}

\end{tabular}}
\begin{minipage}{\textwidth}
\footnotesize
*** $p<0.01$, ** $p<0.05$, * $p<0.1$.\\
\textit{Notes:} The treatment variables are defined using alternative lag windows over the previous three, four, and five years. 
The 4-year-window models, reported in columns (2), (5), and (8), include 761,236 observations; compared with the 3-year-window models, 
13,059 observations are dropped because the required longer lagged treatment information is unavailable. The 5-year-window models, reported 
in columns (3), (6), and (9), include 746,245 observations; compared with the 3-year-window models, 28,050 observations are dropped for the
 same reason. Across specifications, the estimated coefficients are larger in absolute magnitude for the shorter treatment windows for both 
 harvest and thinning. This pattern indicates that the wildfire-mitigation effects of forest treatments are strongest in the years immediately 
 following treatment and gradually diminish as the time since treatment increases.
 Such attenuation is consistent with the gradual reaccumulation of vegetation and combustible fuels following harvest or thinning, 
 which can reduce the effectiveness of earlier treatments in mitigating subsequent wildfire risk.
\end{minipage}
\end{table}

\begin{table}[!htbp]
\centering
\caption*{Table A3. Robustness tests with alternative IV}
\resizebox{\textwidth}{!}{
{
\def\sym#1{\ifmmode^{#1}\else\(^{#1}\)\fi}
\begin{tabular}{l*{6}{c}}
\toprule
                &\multicolumn{3}{c}{Harvest}                             &\multicolumn{3}{c}{Thinning}                            \\\cmidrule(lr){2-4}\cmidrule(lr){5-7}
                &\multicolumn{1}{c}{(1)}&\multicolumn{1}{c}{(2)}&\multicolumn{1}{c}{(3)}&\multicolumn{1}{c}{(4)}&\multicolumn{1}{c}{(5)}&\multicolumn{1}{c}{(6)}\\
                &\multicolumn{1}{c}{State IV}&\multicolumn{1}{c}{State-Eco IV}&\multicolumn{1}{c}{County IV}&\multicolumn{1}{c}{State IV}&\multicolumn{1}{c}{State-Eco IV}&\multicolumn{1}{c}{County IV}\\
\midrule
Harvest         &  -0.4453\sym{***}&  -0.2790\sym{***}&  -0.1950\sym{***}&                  &                  &                  \\
                & (0.0830)         & (0.0685)         & (0.0750)         &                  &                  &                  \\
\addlinespace
Thinning        &                  &                  &                  &  -0.1872\sym{***}&  -0.1831\sym{***}&  -0.1065\sym{**} \\
                &                  &                  &                  & (0.0435)         & (0.0437)         & (0.0530)         \\
\midrule
Kleibergen-Paap Wald F statistic&   214.64         &   262.55         &   155.54         &   488.11         &   415.47         &   165.12         \\
Observations    &  761,236         &  761,236         &  761,236         &  761,236         &  761,236         &  761,236         \\
Plot FE         &\checkmark         &\checkmark         &\checkmark         &\checkmark         &\checkmark         &\checkmark         \\
Year FE         &\checkmark         &\checkmark         &\checkmark         &\checkmark         &\checkmark         &\checkmark         \\
Species group FE&\checkmark         &\checkmark         &\checkmark         &\checkmark         &\checkmark         &\checkmark         \\
Ecoregion FE    &\checkmark         &\checkmark         &\checkmark         &\checkmark         &\checkmark         &\checkmark         \\
\bottomrule
\multicolumn{7}{l}{\footnotesize Standard errors in parentheses}\\
\multicolumn{7}{l}{\footnotesize \sym{*} \(p<0.10\), \sym{**} \(p<0.05\), \sym{***} \(p<0.01\)}\\
\end{tabular}
}

}
\begin{minipage}{\textwidth}
\footnotesize
*** $p<0.01$, ** $p<0.05$, * $p<0.1$.\\
\textit{Notes:} Columns (1)–(3) report IV estimates for harvest using delocalized shift-share instruments constructed at three levels of geographic aggregation: state, state-ecoregion, and county-ecoregion, respectively. Columns (4)–(6) report the corresponding estimates for thinning using the same set of instruments.
\end{minipage}
\end{table}



\begin{table}[!htbp]
\centering
\caption*{Table A4. Plot-level robustness tests without Texas and Oklahoma}
\resizebox{\textwidth}{!}{\begin{tabular}{llll}
\begin{tabular}{lcccc}
\toprule
& \multicolumn{2}{c}{Harvest} & \multicolumn{2}{c}{Thinning} \\
\cmidrule(lr){2-3}\cmidrule(lr){4-5}
& (1) & (2) & (3) & (4)  \\
& \multicolumn{1}{c}{Wildfire event} & \multicolumn{1}{c}{Wildfire event} & \multicolumn{1}{c}{Wildfire event} & \multicolumn{1}{c}{Wildfire event} \\
\midrule
Harvest             &     -0.4453\sym{***}&     -0.4585\sym{***}&                     &                     \\
                    &    (0.0830)         &    (0.0963)         &                     &                     \\
\addlinespace
Thinning            &                     &                     &     -0.1872\sym{***}&      0.0092         \\
                    &                     &                     &    (0.0435)         &    (0.0450)         \\
\midrule
Kleibergen-Paap Wald F statistic&      214.64         &      181.85         &      488.11         &      438.74         \\
Observations        &     761,236         &     542,064         &     761,236         &     542,064         \\
Plot FE             &  \checkmark         &  \checkmark         &  \checkmark         &  \checkmark         \\
Year FE             &  \checkmark         &  \checkmark         &  \checkmark         &  \checkmark         \\
Species group FE    &  \checkmark         &  \checkmark         &  \checkmark         &  \checkmark         \\
Ecoregion FE        &  \checkmark         &  \checkmark         &  \checkmark         &  \checkmark         \\
\bottomrule
\end{tabular}

\end{tabular}}
\begin{minipage}{\textwidth}
\footnotesize
*** $p<0.01$, ** $p<0.05$, * $p<0.1$.\\
\textit{Notes:} Columns (1) and (3) use the full estimation sample. 
Columns (2) and (4) exclude Texas and Oklahoma, where artificially 
regenerated forests are limited, to reduce concerns that the shift-share 
instruments may rely on extrapolation beyond areas with adequate variation 
in artificial regeneration.
\end{minipage}
\end{table}
 



\begin{table}[!htbp]
\centering
\caption*{Table A5. Robustness tests for alternative wildfire indicators}
\resizebox{\textwidth}{!}{%
\begin{tabular}{l*{4}{c}}
\\\midrule
                &\multicolumn{1}{c}{Log wildfire area (MTBS)}&\multicolumn{1}{c}{Wildfire event (MTBS)}&\multicolumn{1}{c}{Percent fire (MTBS)}\\
Panel A. Harvest results \\
& (1) & (2) & (3) \\
\midrule
Log average harvest area&   -0.462\sym{***}&   -0.061\sym{***}&   -0.001\sym{***}\\
                &  (0.114)         &  (0.015)         &  (0.000)         \\
\midrule
Kleibergen-Paap Wald F statistic&    59.61         &    59.61         &    59.61         \\
Observations    &    30312         &    30312         &    30312         \\
County FE       &$\checkmark$         &$\checkmark$         &$\checkmark$         \\
Year FE         &$\checkmark$         &$\checkmark$         &$\checkmark$         \\
\midrule

                &\multicolumn{1}{c}{Log wildfire area (MTBS)}&\multicolumn{1}{c}{Wildfire event (MTBS)}&\multicolumn{1}{c}{Percent fire (MTBS)}\\
Panel B. Thinning results \\
& (4) & (5) & (6) \\
\midrule
Log average thinning area&   -0.211\sym{***}&   -0.029\sym{***}&   -0.000         \\
                &  (0.049)         &  (0.006)         &  (0.000)         \\
\midrule
Kleibergen-Paap Wald F statistic&   137.53         &   137.53         &   137.53         \\
Observations    &    30312         &    30312         &    30312         \\
County FE       &$\checkmark$         &$\checkmark$         &$\checkmark$         \\
Year FE         &$\checkmark$         &$\checkmark$         &$\checkmark$         \\
\midrule

\end{tabular}%
}
\begin{minipage}{\textwidth}
\footnotesize
*** $p<0.01$, ** $p<0.05$, * $p<0.1$.\\
\textit{Notes:} Percent wildfire is defined as the share of forest 
   area burned by wildfire in that county-year. Because these measures normalize by county forest area, they reduce the influence of
    county size and show that the main results, especially the harvest treatment, are not solely driven by larger counties. 
\end{minipage}
\end{table}





\clearpage
\begin{landscape}
\begin{table}[!htbp]
\centering
\caption*{Table A6. Robustness of the Effects of Harvest and Thinning on County-Level Wildfire Burned Area}
\label{tab:A3}

\resizebox{\textwidth}{!}{%
\begin{tabular}{l*{12}{c}}
\\\midrule
                &\multicolumn{6}{c}{Harvest}                                                                                      &\multicolumn{6}{c}{Thinning}                                                                                     \\\cmidrule(lr){2-7}\cmidrule(lr){8-13}
                &\multicolumn{1}{c}{Log wildfire area}&\multicolumn{1}{c}{Log average treatment}&\multicolumn{1}{c}{Log wildfire area}&\multicolumn{1}{c}{Log wildfire area}&\multicolumn{1}{c}{Log wildfire area}&\multicolumn{1}{c}{Log wildfire area}&\multicolumn{1}{c}{Log wildfire area}&\multicolumn{1}{c}{Log average treatment}&\multicolumn{1}{c}{Log wildfire area}&\multicolumn{1}{c}{Log wildfire area}&\multicolumn{1}{c}{Log wildfire area}&\multicolumn{1}{c}{Log wildfire area}\\
& (1) & (2) & (3) & (4) & (5) & (6) & (7) & (8) & (9) & (10) & (11) & (12) \\
& FE & First stage & IV: 4-year & IV: 3-year & IV: 5-year & IV no COVID & FE & First stage & IV: 4-year & IV: 3-year & IV: 5-year & IV no COVID \\
\midrule
Log average treatment&  -0.0036         &                  &  -0.4622\sym{***}&  -0.2916\sym{***}&  -0.7207\sym{***}&  -0.5702\sym{***}&  -0.0093         &                  &  -0.2107\sym{***}&  -0.2147\sym{***}&  -0.1734\sym{**} &  -0.2692\sym{***}\\
                & (0.0042)         &                  & (0.1138)         & (0.0797)         & (0.1696)         & (0.1214)         & (0.0063)         &                  & (0.0485)         & (0.0427)         & (0.0708)         & (0.0515)         \\
Shift-share instrument&                  &   0.0001\sym{***}&                  &                  &                  &                  &                  &   0.0003\sym{***}&                  &                  &                  &                  \\
                &                  & (0.0000)         &                  &                  &                  &                  &                  & (0.0000)         &                  &                  &                  &                  \\
\midrule
Kleibergen-Paap Wald F statistic&        .         &        .         &    59.61         &    57.78         &    42.83         &    60.47         &        .         &        .         &   137.53         &   141.86         &    88.92         &   132.73         \\
Observations    &    30312         &    30312         &    30312         &    30312         &    29049         &    25260         &    30312         &    30312         &    30312         &    30312         &    29049         &    25260         \\
County FE       &$\checkmark$         &$\checkmark$         &$\checkmark$         &$\checkmark$         &$\checkmark$         &$\checkmark$         &$\checkmark$         &$\checkmark$         &$\checkmark$         &$\checkmark$         &$\checkmark$         &$\checkmark$         \\
Year FE         &$\checkmark$         &$\checkmark$         &$\checkmark$         &$\checkmark$         &$\checkmark$         &$\checkmark$         &$\checkmark$         &$\checkmark$         &$\checkmark$         &$\checkmark$         &$\checkmark$         &$\checkmark$         \\
\midrule

\end{tabular}%
}

\vspace{0.5em}

\begin{minipage}{\textwidth}
\footnotesize
*** $p<0.01$, ** $p<0.05$, * $p<0.1$.\\
\textit{Notes:} Columns (1)--(6) report estimates for harvest, and columns (7)--(12) report estimates for thinning. Columns (4) and (10) use a three-year treatment window,
 which reduces the sample by one year. Columns (5) and (11) use a five-year treatment window. Columns (6) and (12) exclude observations from year 2020, which reduces the
  sample by three years.
\end{minipage}

\end{table}
\end{landscape}
\clearpage

\end{document}